\section{Expansion and the sequential probability law}
\label{sec:expansion}

We now choose the permutation tables to make every auxiliary incidence
expand.  The choice must also retain a quantitative probability estimate
for pruning: at one source-block update, conditioning on expansion may
increase the probability of a joint test of $k$ current tables by at most
$e^{\varepsilon_s k}$, where $\varepsilon_s\to0$ as $s\to\infty$.
Proposition~\ref{prop:sequential-law} supplies both conclusions for one
law on the full tower.  The expansion bound concerns the auxiliary
incidences before vertex deletions.

Throughout the finite estimates, fix a family at a marked cutoff, all
its earlier tables, and one extra-bit slice.  Only its current
active-prefix tables are random, and the reference law makes them
independent uniform permutations.  The auxiliary prescription forces
the current gate open.  We first bound the mean incidence operator,
then select tables for which all cuts expand, controlling the change of
law at the same time.  Every bound is uniform in the earlier history
and the height; this uniformity will permit successive choices along
the tower.

Put $\theta=2^{1-b}$.  Choose the fixed integer $b$ sufficiently large
that $\sqrt{5\theta}<3/5$.  All averages on either vertex side are
probability averages.  For one fixed-extra slice, write $A,B$ for its
common-cutoff coordinate spaces and define
\[
 (Nf)(a)=\frac1d\sum_{e:e^-=a}f(e^+),\qquad
 (N^*g)(b')=\frac1D\sum_{e:e^+=b'}g(e^-).
\]
The sums count slots, including parallel edges.  Since $d|A|=D|B|$,
these are adjoint contractions that preserve constants.  Let
$\mathsf{M}=\mathbb E N$, where the expectation is only over the current tables.
The same statements hold for $\mathsf{M}$.

\subsection{The current matching variables}

The triangular coordinates from Section~\ref{sec:labeling} identify
exactly what the current random permutations match.  Slots are the
objects in these matchings; different slots of one physical vertex
remain distinct.

\begin{lemma}[Fresh matching coordinates]\label{lem:fresh-matching-coordinates}
Fix a family occurrence, an extra slice, and every table before its last
common source block.  Under the independent uniform current-table law,
its auxiliary incidence is randomized by independent
perfect matchings between pairs of clusters of $2^t$ slot incidences, one
matching for each full final-prefix table key.  Each possible edge in a
cluster pair is permitted.  No later table affects this incidence.
\end{lemma}

\begin{proof}
Fix the serial number, full label head and preceding raw $u$ blocks.  If
the last common block is B, its remaining B data are exactly the $t$
bits of raw $u_B$ in that block; the known correction merely translates
them.  Its remaining A data are exactly the corresponding $t$ evaluation
bits.  All other coordinates are determined by
Lemma~\ref{lem:incidence-inversion}.  The forced-on current prefix table
is an arbitrary bijection between these two $2^t$-element sets.
Different complete keys specify disjoint clusters of slot incidences.
The A-terminal case reverses the roles.  The padding proof in
Lemma~\ref{lem:incidence-inversion} shows that higher tables never act on
this incidence.
\end{proof}

\subsection{A conditional Fourier estimate}

The coefficient multiplicity bound gives cancellation when we average
the evaluations $V+K_aU$ over slots.  We use the following conditional
form so that earlier raw prefixes and their correction offsets may be
held fixed.

\begin{lemma}\label{lem:conditional-fourier}
Conditional on data $C$, let $U\in\F_2^t$ and $V\in\F_2^w$ be
independent uniform vectors, independent also of auxiliary data $W$.
For $a$ in a set of $q_0$ slots let $K_a:\F_2^t\to\F_2^w$ be linear,
let $o_a=o_a(C,W)$, and let $f_a(C,z)$ have no separate dependence on
$U,V,W$.  Suppose every class of slots with a common value of
$\zeta K_a$, for any nonzero functional $\zeta\in(\F_2^w)^*$,
has at most $\theta q_0$ members.  Set
\[
 F=\frac1{q_0}\sum_a f_a(C,V+K_aU+o_a),\qquad
 \mu(C)=\frac1{q_0}\sum_a2^{-w}\sum_z f_a(C,z).
\]
Then $\mathbb E(F\mid C,W,U)=\mu(C)$ and
\begin{equation}\label{eq:fourier-conditional-bound}
 \mathbb E|F-\mu(C)|^2
 \leq\theta\,\mathbb E\left[\frac1{q_0}\sum_a2^{-w}
                       \sum_z|f_a(C,z)|^2\right].
\end{equation}
\end{lemma}

\begin{proof}
Translation of uniform $V$ proves the mean assertion.  Fix $C,W$ and
write $\widehat f_a(\zeta)=2^{-w}\sum_z f_a(C,z)(-1)^{\zeta z}$.
Orthogonality, first in $V$ and then in $U$, gives
\[
 \mathbb E_{U,V}|F-\mu(C)|^2
 =\frac1{q_0^2}\sum_{\zeta\ne0}\sum_{\mathcal C}
       \left|\sum_{a\in\mathcal C}(-1)^{\zeta o_a}
                          \widehat f_a(\zeta)\right|^2,
\]
where $\mathcal C$ ranges over classes with equal $\zeta K_a$.
Cauchy--Schwarz bounds each square by
$|\mathcal C|\sum_{a\in\mathcal C}|\widehat f_a(\zeta)|^2$.
The multiplicity assumption and Parseval's identity prove the result
after integration over $C,W$.  In particular, integration over $W$
introduces no variance of the conditional mean.
\end{proof}

The coefficient pools of Section~\ref{sec:labeling} give this
multiplicity bound on every full block and every active initial part of
length at least $b$.  The first $A$ block has length $r$; all subsequent
blocks have length $s$.  The divisibility choices make every nonempty
terminal initial part have length at least $b$.

We record why the right side of
\eqref{eq:fourier-conditional-bound} is bounded by the original input
norm in its applications.  A uniform vertex with a uniform incident
slot has a uniform neighbor, counted with the appropriate multiplicity.
Consequently the slot average of squared neighbor values is exactly the
squared input norm.  Averaging any still random table values can only
decrease this quantity, by Jensen's inequality.  Finally $V$ makes the
evaluation argument uniform conditional on all other data.  Thus the
individual Fourier norms in the Lemma are precisely bounded by these
slot norms; no factor depending on the number of blocks is introduced.

\subsection{Two interacting prefix bands}

To apply the Fourier estimate to the whole incidence, decompose a
function into orthogonal increments as longer state prefixes are revealed.
Triangular inversion couples only two neighboring levels of these
prefix decompositions.  The only additional estimate is at the terminal
block, where the current prefix permutation supplies the missing
average.

Let $P_h^A$ be expectation onto the first $r+hs$ bits of both raw
$A$ coordinates, and let $P_j^B$ be expectation onto the first $js$
bits of both raw $B$ coordinates.  Padded coordinates are deterministic
zero.  The filtrations therefore terminate at the identities.  In
particular $P_0^B=C_B$, the projection onto constants.  Write $C_A$
for the corresponding projection on $A$.

For a fixed $A$ slot, a function in $\operatorname{ran}P_j^B$ can
be evaluated at its neighbor using only
\begin{equation}\label{eq:forward-prefix-data}
 u_A[0,r+(j-1)s),\qquad Q_A[0,r+js).
\end{equation}
The first $r$ evaluation bits give the label head.  Successive inversion
of the triangular permutations then gives $u_B$ through its $j$th
block; the equation for $Q_B$ uses $u_A$ only through its $(j-1)$st
block to give $v_B$ through length $js$.
In the reverse direction, a fixed $B$ slot evaluates a function in
$\operatorname{ran}P_h^A$ using only
\begin{equation}\label{eq:reverse-prefix-data}
 u_B[0,hs),\qquad Q_B[0,(h+1)s).
\end{equation}
Here the bucket is already specified by the slot.  The remaining head
and then $u_A$ are recovered from the evaluation bits, and the equation
for $Q_A$ gives the required $v_A$ bits.

These descriptions also apply to the corrected coordinates.  Their
extra bits are fixed in the slice.  Every correction gate ends before
both coupled blocks, so its offset is determined by the preceding raw
prefix.  The same is true of an ordinary permutation gate.  A current
gate is forced open, hence constant.  No unknown current raw input
enters a key for its own source block.  The identity suffix at a mark
and the prescribed subsequent identity blocks make the assertions
valid beyond an actual cutoff as well.  These are the reasons for
using raw prefixes in the keys and the lag specified in
Section~\ref{sec:labeling}.

Since multiplication is block diagonal, there is no spill into a
following evaluation block.  Equations
\eqref{eq:forward-prefix-data}--\eqref{eq:reverse-prefix-data} imply
\begin{equation}\label{eq:two-prefix-bands}
 \operatorname{ran}(\mathsf{M}P_j^B)\subseteq\operatorname{ran}P_j^A,
 \qquad
 \operatorname{ran}(\mathsf{M}^*P_h^A)\subseteq\operatorname{ran}P_{h+1}^B.
\end{equation}
They hold before as well as after averaging current tables.
For every full interval that is active, Lemma~\ref{lem:conditional-fourier}
also gives
\begin{align}
 \|(1-P_{j-1}^A)\mathsf{M}P_j^B\|^2&\leq\theta,
       \label{eq:forward-full}\\
 \|(1-P_h^B)\mathsf{M}^*P_h^A\|^2&\leq\theta.
       \label{eq:reverse-full}
\end{align}
For \eqref{eq:forward-full}, condition on the preceding $A$ prefix.
The new data in \eqref{eq:forward-prefix-data} occur solely as
$V+K_aU+o_a$, with $U$ the new $A$ input block and $V$ its evaluation
block of $v_A$ bits.  The direct input in that equation stops before
$U$.  The $A$ coefficient does not depend on the unknown bucket in
the head.  For \eqref{eq:reverse-full}, condition on $P_h^B$ and use
the next $B$ input and evaluation blocks.  This time the direct input
in \eqref{eq:reverse-prefix-data} again stops before $U$; the slot
already fixes the bucket, and its coefficient does not depend on the
unknown suffix of the head.  Thus neither application has a separate
dependence on the fresh $U$.  The slot-norm argument above completes
both estimates, including $h=0$.

\begin{lemma}[The conditional mean]\label{lem:conditional-mean}
For every fixed earlier-table history, every marked cutoff, and every
extra-bit slice,
\[
 \|\mathsf{M}:1_B^\perp\longrightarrow1_A^\perp\|
       \leq\eta:=\sqrt{5\theta}<3/5.
\]
The same bound holds for the complete adjacent-height incidence.
\end{lemma}

\begin{proof}
\emph{The full blocks.}
Use the orthogonal difference projections
\[
 E_0^A=P_0^A-C_A,\quad E_h^A=P_h^A-P_{h-1}^A\ (h\geq1),
 \qquad E_j^B=P_j^B-P_{j-1}^B\ (j\geq1).
\]
Equation~\eqref{eq:two-prefix-bands} implies
$E_h^A\mathsf{M}E_j^B=0$ unless $h\leq j\leq h+1$.
The band $h=j$ is bounded by \eqref{eq:forward-full}, and the band
$j=h+1$ by the adjoint of \eqref{eq:reverse-full}, whenever the
blocks in question are full.  Within either band the input spaces
are mutually orthogonal, as are the output spaces.  Each band
therefore has norm at most $\sqrt\theta$, and their sum at most
$2\sqrt\theta$.

\emph{The terminal block.}
There are two possible terminal forms.  If the final input block is
a $B$ block, write
\[
 \ell_A=r+ns,\quad \ell_B=ns+t,\qquad
 |v_A|=r+ns+t,\quad |v_B|=(n+1)s,\quad b\leq t\leq s.
\]
For the input differences with $1\leq j\leq n$, all blocks required
in both bands are full, and consequently
\begin{equation}\label{eq:low-B-mean}
 \|\mathsf{M}(P_n^B-C_B)\|\leq2\sqrt\theta.
\end{equation}
Fix a $B$ slot.  Inversion through the preceding blocks determines
all of $u_A$ and all but the last $t$ bits of $v_A$ using
$u_B[0,ns)$ and $Q_B[0,(n+1)s)$.  The current prefix permutation
uniformly averages the remaining $t$ bits of $Q_A$.  Its coefficient
term $K_Au_A$ ends at $r+ns$, so these are exactly the remaining
$v_A$ bits, up to a known correction.  The matching key contains
only the label and preceding raw inputs.  After this average, its
slot value has no separate dependence on the last $t$ raw $u_B$
bits.  Conditional on $P_n^B$, it is therefore a function of
$V+K_aU+o_a$, with $U\in\F_2^t$ and $V\in\F_2^s$.
Lemma~\ref{lem:conditional-fourier} gives
\[
 \|(1-P_n^B)\mathsf{M}^*\|^2\leq\theta.
\]
Combining this orthogonally with the adjoint of
\eqref{eq:low-B-mean} gives $\|\mathsf{M}^*g\|^2\leq5\theta\|g\|^2$
for centered $g$.

If the final input block is an $A$ block, the lengths are instead
\[
 \ell_B=(n+1)s,\quad \ell_A=r+ns+t,\qquad
 |v_A|=r+(n+1)s,\quad |v_B|=(n+1)s+t.
\]
The full-block argument on $E_h^A$, $0\leq h\leq n$, gives
$\|\mathsf{M}^*(P_n^A-C_A)\|\leq2\sqrt\theta$.
The current permutation averages the final $t$ evaluation bits of
$Q_B$, on which $K_Bu_B$ is zero.  Conditional on $P_n^A$, the
remaining slot function depends on the final $t$ input bits of $A$
only through their full $s$-bit evaluation.  Hence
$\|(1-P_n^A)\mathsf{M}\|^2\leq\theta$, and the same orthogonal combination
proves the assertion.  Both arguments include $t=s$; the initial
$r$-block is never terminal because the tower starts after many
complete pairs.

\emph{Combining the extra-bit slices.}
If the larger side
is $A$, its space is $A_0\times\F_2^b$ and the complete operator
has components $\mathsf{M}_e$.  For centered $f$ on $B$,
$\|\mathsf{M}f\|^2=\mathbb E_e\|\mathsf{M}_ef\|^2\leq\eta^2\|f\|^2$.
If the larger side is $B$, use the adjoint version of this argument.
Equivalently, write an input as $f_e=c_e+f_e^0$; its overall mean
zero says $\mathbb E_ec_e=0$, so the constants cancel before
Jensen's inequality is applied to $\mathbb E_e\mathsf{M}_ef_e^0$.
\end{proof}

\subsection{Cut concentration}

The mean estimate gives a lower bound for every expected cut.  The
matching coordinates now let us turn this into a failure-probability
bound for each fixed cut.

We now condition only on earlier tables and one extra-bit slice.
Lemma~\ref{lem:fresh-matching-coordinates} identifies each current
table with a uniform bijection between two fibers of slots of size
$2^t$.  Its key fixes all other endpoint coordinates.  Distinct keys
give disjoint such fibers of slots, and the current tables are
independent.  A physical vertex may occur in several fibers through
its different slots; this does not identify their permutation
variables.

Write $\deg(v)=d$ on $A$ and $D$ on $B$, and
$\vol(S)=\sum_{v\in S}\deg(v)$.  The total volume is $V_0=2d|A|$.
Let $c(S)$ count edges crossing from $S$ to its complement, with
multiplicity.  The mean random-walk operator on the union of the
two sides is the self-adjoint block operator with entries $\mathsf{M},\mathsf{M}^*$.
On the orthogonal complement of global constants its spectrum has
upper bound $\eta$; its additional bipartition eigenvector has
eigenvalue $-1$.  Applying its quadratic form to the centered
indicator of $S$ gives
\begin{equation}\label{eq:expected-cut}
 \mathbb E c(S)\geq(1-\eta)\vol(S)
                  \left(1-\frac{\vol(S)}{V_0}\right)
               \geq\frac15\vol(S)
\end{equation}
whenever $\vol(S)\leq V_0/2$.

For a current table $a$, let $Z_a$ count its matched pairs whose two
endpoints belong to $S$.  If $n_a$ is the fiber size, $u_a$ the
number of its $A$ slots in $S$, and $v_a$ the corresponding number
of $B$ slots, then $Z_a$ is hypergeometric with these parameters.
Moreover
\[
 c(S)=\vol(S)-2\sum_a Z_a,\qquad
 \sum_a u_a\leq\vol(S).
\]
For completeness, sampling $u_a$ distinct entries from a population
of $v_a$ ones and $n_a-v_a$ zeros gives, for every real $\lambda$,
\[
 \mathbb E e^{\lambda(Z_a-\mathbb EZ_a)}
       \leq e^{\lambda^2u_a/8}.
\]
Indeed the average of products of $u_a$ distinct positive numbers
is at most the $u_a$th power of their arithmetic mean: this follows
by repeatedly replacing two unequal numbers by their average in
the elementary symmetric polynomial, which increases it.  Apply
this to the numbers $1,e^\lambda$.  The resulting bound is the
moment generating function for sampling with replacement.  For
one Bernoulli variable its centered log moment generating function
has value and first derivative zero at zero and second derivative
at most $1/4$ everywhere, proving the displayed bound.
Independence and exponential Markov's inequality therefore yield
\begin{equation}\label{eq:cut-failure}
 \mathbb P\{c(S)<.05\vol(S)\}
 \leq \exp\{- .01125\vol(S)\}
 \leq \exp\{-c_0d_{\min}|S|\},\qquad c_0=10^{-3},
\end{equation}
where $d_{\min}=\min(d,D)$.  The numerical constant uses
$\sum_aZ_a-\mathbb E\sum_aZ_a>.075\vol(S)$ and the choice
$\lambda=.3$.

\subsection{The finite selection lemma and its conditioning bound}

Join two physical vertices when some choice of a current table
could join them.  This \emph{possible graph} has maximum degree
$\Delta\leq2^{C s}$, where $C$ may depend on the fixed number of
types and on $b$.  Each current table meets at most $2^{C s}$
vertices; the number of tables meeting a vertex has the same bound.
These bounds do not depend on the height.  It suffices to avoid
the cut failures in \eqref{eq:cut-failure} for nonempty connected
sets $S$ in the possible graph satisfying $\vol(S)\le V_0/2$.
In fact any failed cut of at most half the volume has a failed
connected component, since both volume and crossing add over
those components.

Let $A_S$ be such a cut-failure event.  It depends only on tables
meeting $S$.  Two events with no common table are independent,
jointly with any collection of other events using disjoint tables.
If their table sets meet, some vertex of one is within distance
two of a vertex of the other in the possible graph.  A connected
$j$-set containing a specified vertex is encoded by a walk of
length $2(j-1)$ obtained from a deterministically chosen rooted
spanning tree.  Thus the number of neighboring $j$-set events is
at most
\begin{equation}\label{eq:dependency-count}
 |S|\Delta^{2j+C_1}
\end{equation}
for an absolute enlargement $C_1$.

Assign weights $x_S=\exp(-c_0d_{\min}|S|/2)$.  Equations
\eqref{eq:dependency-count} and the exponential growth of
$d_{\min}=2^{.40005s}$ imply, uniformly in $S$,
\[
 \sum_{T\sim S}x_T
 \leq |S|\sum_{j\geq1}\Delta^{2j+C_1}e^{-c_0d_{\min}j/2}
       =|S|\delta_s,\qquad\delta_s\longrightarrow0.
\]
For sufficiently large $s$, all weights are below $1/2$ and
$2\delta_s\leq c_0d_{\min}/2$.  Consequently
\begin{equation}\label{eq:lll-criterion}
 \mathbb P(A_S)\leq x_S\prod_{T\sim S}(1-x_T).
\end{equation}

This is the finite asymmetric Lov\'asz local lemma in its product form
\cite[Theorem~1.1]{MT10}; the original symmetric lemma is due to
Erd\H{o}s and Lov\'asz \cite[\S2]{EL75}. We include the conditional-probability
induction because its quantitative conditioning bound is also needed.  For a finite family
satisfying \eqref{eq:lll-criterion}, induction simultaneously proves
positivity of avoidance probabilities and the following conditional
bound.  At a given size, positivity follows from the smaller-size
bounds by the product rule; the empty avoidance event has probability
one.  For any $A_i$ the bound is
\[
 \mathbb P\left(A_i\,\middle|\,\bigcap_{j\in J}A_j^c\right)
       \leq x_i\qquad(i\notin J).
\]
Split $J$ into the neighbors $J_1$ of $i$ and its remaining set
$J_2$.  The numerator after conditioning on avoidance in $J_2$
is at most $\mathbb P(A_i)$, by independence.  If $J_1$ is empty,
independence proves the desired bound immediately.  Otherwise the denominator
for also avoiding $J_1$ is, by the product rule and the induction
hypothesis, at least $\prod_{j\in J_1}(1-x_j)$.  The criterion
now proves the induction step.  The same product rule proves
that avoidance of every event has strictly positive probability.

We also need the standard conditional-distribution estimate for an
additional event; see \cite[Theorem~2.1]{HSS11}.  Let $B$ involve
a specified set of $k$ current tables.  Let $J_1$ be the cut events
sharing one of these tables and $J_2$ all the others.  Independence
gives
\[
 \mathbb P\left(B\,\middle|\,\bigcap_S A_S^c\right)
 \leq\mathbb P(B)\prod_{S\in J_1}(1-x_S)^{-1}.
\]
For each specified table, its at most $2^{Cs}$ vertices and the
connected-set count give
\[
 \sum_{S\in J_1}x_S
 \leq k\,2^{Cs}\sum_{j\geq1}\Delta^{2j+C_1}
                          e^{-c_0d_{\min}j/2}.
\]
It follows that, for a number $\varepsilon_s\to0$ independent of
height and earlier history,
\begin{equation}\label{eq:current-event-inflation}
 \mathbb P\left(B\,\middle|\,\text{all current cuts succeed}\right)
       \leq e^{\varepsilon_s k}\mathbb P(B).
\end{equation}
The same inequality applies to any nonnegative function of those
$k$ tables, by integration of its level sets.  It concerns a joint
test of current variables.  It does not assert a bound for the
next query after conditioning on arbitrary earlier answers from
the same block.

\subsection{A gap in every chosen incidence}

The selected cut bounds give a spectral gap by the degree-weighted
Cheeger argument; compare \cite[Lemma~3.3]{SJ89} and
Chung~\cite[Section~3, Theorem~1]{ChungCheeger}.  We include the
argument in the bipartite form needed here, keeping its normalization
and edge multiplicities.
If all the selected cuts satisfy $c(S)\geq h_0\vol(S)$ with
$h_0=.05$, the ordinary weighted spectral gap is at least
$h_0^2/2$.
For a real function $f$ of degree-weighted mean zero, choose a
weighted median $m$ and put $g_+=(f-m)_+$ and $g_-=(m-f)_+$.
Each has support of volume at most half.  Integration over the
level sets of $g^2$, for either one, yields
\[
 \sum_{\{u,v\}}|g(u)^2-g(v)^2|
       \geq h_0\sum_v\deg(v)g(v)^2.
\]
Cauchy--Schwarz and
$\sum_{\{u,v\}}(g(u)+g(v))^2\leq2\sum_v\deg(v)g(v)^2$
give
\[
 \sum_{\{u,v\}}(g(u)-g(v))^2
       \geq\frac{h_0^2}{2}\sum_v\deg(v)g(v)^2.
\]
The energy of $f$ is at least the sum of the energies of $g_+$
and $g_-$, while
$\sum_v\deg(v)(f(v)-m)^2\geq\sum_v\deg(v)f(v)^2$.
This proves the claimed gap.  Apply the argument to real and
imaginary parts for complex functions.  Bipartite symmetry sends
an eigenvalue to its negative; the only exceptional modes are
the global constant and bipartition functions.  Hence the
normalized incidence has centered singular norm at most
\begin{equation}\label{eq:uniform-incidence-gap}
 \rho=1-h_0^2/2<1.
\end{equation}
The slice-to-whole argument in Lemma~\ref{lem:conditional-mean}
also transfers this bound to adjacent-height incidences.

\subsection{Choosing the tables along the tower}

At one marked cutoff we have now obtained both a nonempty set of
expanding table choices and the bound
\eqref{eq:current-event-inflation} on conditioning to that set.  Future
tables do not affect this incidence.  We can therefore apply the same
finite choice successively at all marked cutoffs.

\begin{proposition}[Sequential selection]\label{prop:sequential-law}
There is a probability law on all the permutation tables such that
every auxiliary incidence has centered singular norm at most the
fixed number $\rho<1$ in \eqref{eq:uniform-incidence-gap}.
At each individual source-block update, conditional on all earlier
updates, any joint test of $k$ specified current tables has its
probability increased by at most $e^{\varepsilon_s k}$ relative
to independent uniform current tables.  The bound is uniform in
the earlier history and the height, with $\varepsilon_s\to0$.
\end{proposition}

\begin{proof}
Within each of the finitely many families order the source blocks
as $A_0,B_1,A_1,B_2,\ldots$.  Interleave the families by increasing
source-block index and a fixed family order.  At an unmarked
update choose all its active permutation variables uniformly.
At a marked update use the finite conditional distribution given
avoidance of its current cut failures, separately in the extra-bit
slices.  Conditional on all earlier source-block indices, the
current kernels in different families are independent, since a
family's cut conditions use only its own tables.  Thus grouping
the families into one source-block update gives the same bound,
with $k$ the total number of tested tables.  The mark separation
ensures at most one actual cutoff
per family at an update.  Lemma~\ref{lem:conditional-mean} and
\eqref{eq:lll-criterion} show that this distribution exists for
every possible preceding history, and
\eqref{eq:current-event-inflation} proves its test bound.  Blocks
prescribed to be identities have no random variables.

Each update has finitely many possible outcomes.  An independent
uniform number in $[0,1]$ can choose among them by consecutive
intervals of their conditional probabilities.  Iterating on the
countable sequence of updates constructs the asserted law.  A
future update does not change an earlier incidence: the prefix
and padding properties of Section~\ref{sec:labeling} guarantee
this.  Therefore all earlier expansion conditions persist.

When a test whose current-table choices and input arguments are
measurable from earlier updates is evaluated, all
later conditional distributions integrate to one.  Thus its
probability is governed by the current kernel just described,
without an additional conditioning on future expansion successes.
This observation is the form of sequential conditioning used
in the pruning argument below.
\end{proof}
